\section{Discussion: Claim Boundaries, Verification Conditions, and Future Work}

\subsection{Logical separation of soundness and completeness}
Theorem \ref{thm:type-safety-offcritical} provides soundness: under this paper's typed read semantics, any \emph{already addressable and auditable} zero certificate must satisfy critical line constraints. In contrast, completeness ("zero events must be addressable" i.e., read non-$\NULL$) does not automatically follow from null-value semantics and must appear as a bridging assumption (e.g., Assumption \ref{ass:xi-address-completeness}) or be proven via additional construction. This separation ensures reasoning dependency chains are auditable: all inferences involving classical RH must explicitly annotate their bridging premises (e.g., Assumption \ref{ass:bridge-rh}).\par

\subsection{Conditional boundaries and verifiable conditions of BridgeRH}
The role of BridgeRH is to organize "classical zero events" as auditable reads in the auditable object domain; thus its core is not to provide analytic number-theoretic proofs at the same level as the main text, but to transcode difficulty into a group of line-by-line auditable sub-conditions. Its core requirements can be summarized as two boundary constraints:
\begin{enumerate}
  \item \textbf{Object alignment:} establish analytic rigidity consistent with classical object at kernel/determinant level (same zero set and multiplicity), decomposing this consistency into auditable sub-conditions;
  \item \textbf{Event reduction and auditing lock:} reduce classical zero events to PW-addressable addresses, ensuring auditing paths are locked to unitary slice auditing mechanism, making "zero/non-zero/certificate missing ($\NULL$)" offline auditable facts.
\end{enumerate}
The verifiable condition list of BridgeRH, controlled examples, and auditing vocabulary appear in Appendix \ref{app:bridge-rh-obligations}; unified coding directly related to certificate object formats and unitary slice numerical certificate construction appear in Section \ref{sec:conclusion-cluster} and Appendix \ref{app:inject-code-audit}.\par

\subsection{Scope and implementation meaning of three-register closure}
The three-register (scale, phase, prime) closure serves two functions in this paper: first, providing canonical representatives of addressable objects and residual carrier axes for reversible replay; second, providing certificatable continuous visible parameter axes (phase) for unitary slice auditing mechanism and implementing type-safety rejection chains. It must be emphasized: three-register closure is \emph{one of the minimal implementations of the certificate type system}, not an additional ontological assumption on external mathematical objects; within this closure, any assertion depending on radial degrees of freedom ($\rho=\log|u|$) without additional certificate channel degenerates to $\NULL$ by definition (see Theorem \ref{thm:cc-fourth-register-certs} in Section \ref{sec:conclusion-cluster}).\par

At the implementation level, the constraint "visible parameter axis is only phase" can be realized via explicit conjugacy compactification internalization: addition/translation on the real axis can be written as a closed-form Möbius action on unit circle phase domain; when further restricting phase to rational points, exact coding and height budget provide repeatable audit discrete complexity ledger (see Appendix \ref{app:phase-compactification-register-unification}).\par

\subsection{Limitations, falsifiable points, and consistency checks}
This paper's "publishable claims" are bounded by offline-auditable certificate objects, so limitations and falsifiable points can be directly expressed as auditable conditions at the verification rule level:
\begin{itemize}
  \item \textbf{Verification rule dependence:} numerical auditing conclusions depend on fixed outward-rounding interval arithmetic and remainder bound certificates; if rounding model is changed or platform-dependent floating-point non-determinism introduced, the corresponding assertions no longer belong to this paper's auditable object domain (Appendices \ref{app:audit}, \ref{app:inject-code-audit}).\par
  \item \textbf{Repeatability of $\NULL$:} under fixed verification rules and parameters, $\NULL$ judgment of same input should be repeatable; if implementation-dependent inconsistency appears, verifier or certificate format fails auditability premises (Theorem \ref{thm:cc-null-repro} in Section \ref{sec:conclusion-cluster}).\par
  \item \textbf{Replay consistency:} lossless lifting and reversible replay claims require providing decidable $\mathrm{code}/\mathrm{decode}$ and replayer; if there exists input where replayer cannot recover original microstate, the corresponding "extended preservation" claim is directly falsified (Appendices \ref{app:prime-uplift}, \ref{app:inject-code-audit}).\par
  \item \textbf{Local counter-examples of bridging conditions:} each condition of BridgeRH should allow outputting local failure witness; auditable violation of any single condition suffices to negate the corresponding bridging inference chain (Appendix \ref{app:bridge-rh-obligations}).\par
\end{itemize}

\subsection{Future work: paths of formalization and certificatization}
Under the fixed stratification (soundness proven, completeness requires bridging or construction), future work can proceed via "verifiable conditions":
\begin{enumerate}
  \item \textbf{Certificate verifiers and primitive libraries:} implement certificate-carrying modules for interval arithmetic, remainder bounds, determinant/spectral bounds, and contraction criteria, making $\mathsf{Cert}_\exists/\mathsf{Cert}_{\NULL}$ independently re-computable and auditable (Appendix \ref{app:inject-code-audit}).\par
  \item \textbf{Item-by-item completion of BridgeRH:} formalize alignment of kernel objects with classical objects, reduction-injection of zero events, and auditing path lock as three classes of conditions, providing certificate materials for each (Appendix \ref{app:bridge-rh-obligations}).\par
  \item \textbf{Fourth register extension and closure certificates:} if the goal is auditable forms of global completeness, provide certificate carriers for radial degrees of freedom (fourth register) and implement auditable residual certificates for corresponding closure constraints.\par
  \item \textbf{Local--global and gluing obstacles:} adelic local--global closure, sheaf-theoretic gluing obstacle semantics of $\NULL$, and defect witness extraction constructions, as candidate structural-level extensions of "failure attribution formats", appear in Appendix \ref{app:adelic-sheaf-null-witness}.\par
\end{enumerate}

To match these advances, two classes of offline-auditable channels have been provided unified exposition in appendices: one transcodes $H^2$-gluing obstacles into auditable invariants via spectral flow winding numbers (Appendix \ref{app:spectral-flow-cohomology-min-register}); the other provides falsifiable structural lower bounds and refutation via endpoint tail laws and finite-core pole lattice rigidity (Appendix \ref{app:obstruction-tail-pole-rigidity}).\par
